% `mu` is not a length until a *style* says so, and §703 is where the style
% comes in:
%
%     if cur_style<script_style then cur_size:=text_size
%     else cur_size:=script_size*((cur_style-text_style) div 2);
%     cur_mu:=x_over_n(math_quad(cur_size),18);
%
% §730 then converts while *walking* the list, not while reading it:
%
%     glue_node: begin if subtype(q)=mu_glue then
%         begin x:=glue_ptr(q);
%         y:=math_glue(x,cur_mu); delete_glue_ref(x); glue_ptr(q):=y;
%         subtype(q):=normal;
%         end
%     …
%     kern_node: math_kern(q,cur_mu);
%
% So the same `\mkern9mu` is one width in a formula and another in a
% superscript, and at the moment `\mkern` is *read* the style is not yet known
% — a superscript's list is scanned before anything decides it is a script.
% Converting eagerly with the text quad made every `\mkern` under a script one
% size too wide. In this document `\fontdimen6\textfont2` is 10.95003pt and
% `\scriptfont2`'s is 8.50012pt, so `\mkern9mu` is 5.4749pt against 4.25006pt —
% and amsmath builds every `\xrightarrow` as
% `$\scriptstyle\mkern5mu{…}\mkern9mu$`, so every extensible arrow was 2.4pt
% too wide. At that width one more `\cleaders` copy of the arrow's shaft fitted
% than TeX had room for, and the arrow came out a stroke longer.
%
% The inter-atom skips were already right — `\thinmuskip` and friends are
% converted where the style is known — so `$\scriptstyle x+x$` never showed it.
% Only `\mkern` and `\mskip` did.
%
% And §720's "simplify a trivial box", which is the other half of what a clean
% box holds:
%
%     q:=list_ptr(x);
%     if is_char_node(q) then
%       begin r:=link(q);
%       if r<>null then if link(r)=null then
%         if not is_char_node(r) then if type(r)=kern_node then
%           begin free_node(r,small_node_size); link(q):=null;
%           end;
%       end
%
% One character followed by exactly one kern: the kern goes. The order is the
% whole of it — this runs *after* `hpack`, so the box keeps the width the
% italic correction gave it and only the node is dropped. `$\scriptstyle y$` is
% 4.30675pt in pdfTeX, glyph plus correction, *and* its box holds one
% character; both are true at once.
\input plain
\showboxdepth=6 \showboxbreadth=60 \tracingonline=0
\def\M#1{\setbox0=\hbox{#1}\showbox0 }
% `\mkern` at each of the three styles, and under a script that has no style
% node at all — which is the case an eager conversion cannot get right.
\M{$x\mkern9mu x$}
\M{$\scriptstyle x\mkern9mu x$}
\M{$\scriptscriptstyle x\mkern9mu x$}
\M{$x^{a\mkern9mu b}$}
\M{$x^{x^{a\mkern9mu b}}$}
\M{$\displaystyle x\mkern9mu x$}
% `\mskip`, which is the same rule for glue, stretch and shrink included.
\M{$\scriptstyle x\mskip9mu plus 3mu minus 1mu x$}
% A `fil` stretch is *not* converted: §716 scales only a finite order.
\M{$\scriptstyle x\mskip9mu plus 1fil x$}
% §720: a subscript's clean box is one character and keeps no kern, while the
% width it was packed to still counts the italic correction.
\M{$p_q$}
\M{$\scriptstyle p_q$}
% Two characters is not a trivial box, so nothing is dropped.
\M{$x_{pq}$}
% §715 puts the removed correction back as a kern when a one-character
% clean box is reboxed to the wider side of a fraction. Otherwise the stretch
% glue centres the bare glyph rather than the width TeX assigned to its box.
\M{${y\over MM}$}
\end
